\section{应用场景与需求分析}
\subsection{用户与使用方式}
主要使用者包括课程学习群中的学生、活动组织者、讨论组成员，以及维护智能体的群管理员。群成员通过自然语言提问；明确提及 Bot 时进入直接请求流程，开启自动参与后，连续群聊也可以构成参与机会。管理员负责授予群服务权限、配置参与节奏，并在控制台检查运行情况。

\begin{table}[htbp]\centering
\caption{校园与群体协作场景}
\begin{tabularx}{\textwidth}{@{}L{22mm}Y Y@{}}\toprule
场景 & 典型问题 & 系统应完成的工作\\\midrule
教学支持 & “数学分析 B1 有哪些教师，评论分别怎样？” & 提取课程线索，查询课程与教师，区分评论样本和客观开课信息，保留来源。\\
活动参与 & “二课最近有哪些学术活动，最多三项？” & 识别类别、时间和数量限制，返回活动名称、时间、状态与详情入口。\\
学业规划 & “2026 级计算机普通主修需要多少学分？” & 指定年级、专业和方案类型；返回相应版本资料，说明研究快照口径。\\
校园生活 & “工作日东区到西区上午有哪些校车？” & 解析方向、日期类型、时段；区分始发与途经时间，注明时刻表版本。\\
服务与管理 & “最近三条校园通知是什么？” & 搜索通知，呈现标题、发布日期和原始链接；必要时查询截止日期。\\
群体协作 & “周五改成周六了，最终什么时候开会？” & 使用同群近期历史识别更正关系，并针对当前问题给出明确结论。\\\bottomrule
\end{tabularx}
\end{table}

科研讨论群可以使用一般问答、分工草案和讨论整理能力；当前材料尚未提供专业科研检索、实验管理或科研成果提升的专项验证，因此作品价值主要落在教学支持、校园生活和群体协作。

\subsection{功能需求}
一次校园查询应能从自然语言定位业务服务，绑定必要参数，读取有来源的资料并形成可理解的回答。信息不足时需要澄清；远端服务失败时应显示失败原因或可用的降级说明，不能把未取得的资料补写为事实。

一次群聊参与应保留发言者、引用关系和当前问题，受到本群配置约束。自动启用只读插件需要已有管理员授权，不能因发现一个新话题就安装任意代码或扩大服务范围。

一次管理操作应让用户看到配置目标和实际状态。QQ 连接、账号登录、Agent 可用、模型配置和 MCP 服务连接应分别表达，因为其中任一项可独立失效。

\subsection{质量需求与可检查目标}
\begin{tabularx}{\textwidth}{@{}L{24mm}Y@{}}\toprule
质量维度 & 可检查目标\\\midrule
语义与事实 & 更正优先于旧安排；课程与教师不混淆；结果中的年份、来源和数量限制保持一致。\\
参与节奏 & 自动参与受概率、冷却和小时上限约束；被关闭的群或能力不能绕过配置执行。\\
可维护性 & 核心包不依赖 AstrBot 的事件对象；新增服务主要修改能力定义、映射和 Provider。\\
可诊断性 & 每次运行能区分模型失败、工具失败、未参与和已生成但未投递等状态。\\
可部署性 & 提供依赖锁文件、容器编排和配置示例；离线流程可在无模型 Key、无 QQ 登录时运行。\\\bottomrule
\end{tabularx}

这些目标用于组织实现与验证，尚未形成真实群长期使用的服务等级承诺。当前没有足够样本支持总体准确率、长期打扰率或 P95 延迟指标。
\source{[S1, S8, S10, S11]。场景示例为说明输入，不是未经执行的效果承诺。}
